\chapter{Conclusiones, competencias y trabajo futuro}
\label{cap:conclusiones}

El Capítulo~\ref{cap:introduccion} planteó el problema de partida y un
objetivo general acompañado de nueve objetivos específicos. Los capítulos
intermedios han desplegado, capa a capa, la plataforma que responde a ese
planteamiento, y el Capítulo~\ref{cap:visualizacion} ha aportado evidencia
reproducible de su comportamiento. Este capítulo cierra el documento
retomando esos objetivos, relacionando el trabajo con los resultados de
aprendizaje del máster, y señalando las limitaciones y el trabajo futuro
que se derivan de todo lo anterior.

\section{Resumen del trabajo}
\label{sec:conclusiones_resumen}

\textbf{Problema y motivación.} La información curricular y de trayectoria
investigadora está dispersa entre documentos individuales, como CVN, y
registros públicos a escala masiva, como ORCID, sin una plataforma
abierta y autoalojada que combine ambas naturalezas de dato, tal y como se
planteó en el Capítulo~\ref{cap:introduccion}.

\textbf{Arquitectura desplegada.} Este trabajo construye una plataforma
\textit{lakehouse}, justificada en el Capítulo~\ref{cap:estado_arte} y
desplegada en el Capítulo~\ref{cap:infraestructura}: Kubernetes, MinIO con
Iceberg, Spark orquestado por Airflow y PostgreSQL. El
Capítulo~\ref{cap:ingesta} ingiere y fusiona ORCID real y CVN sintético, y
el Capítulo~\ref{cap:procesamiento} resuelve su identidad y calcula los
indicadores.

\textbf{Resultados.} El Capítulo~\ref{cap:visualizacion} publica esos
indicadores en un panel verificado en directo, mide el efecto del
paralelismo sobre el rendimiento, y documenta fallos operativos reales
corregidos, incluida una reconstrucción completa desde cero.

\section{Cumplimiento de los objetivos}
\label{sec:conclusiones_objetivos}

El objetivo general planteado en el Capítulo~\ref{cap:introduccion} era
diseñar, desplegar y evaluar una plataforma que integrase CVN sintético y
ORCID real, los procesase con resolución de identidad entre fuentes, y
publicase indicadores en un panel evaluado mediante un banco de pruebas de
rendimiento. Se considera cumplido. La Tabla~\ref{tab:tfm_conclusiones_objetivos}
detalla cada objetivo específico, en el mismo orden del
Capítulo~\ref{cap:introduccion}.

Uno de los objetivos se cumple con garantías parciales, de forma medida y
no oculta: la integración de ORCID con CVN alcanza una precisión del
95,8\% y una cobertura del 74,2\% en su regla de repliegue por nombre y
afiliación, por debajo del objetivo de precisión fijado antes de medir. El
Capítulo~\ref{cap:procesamiento} explicó por qué perseguirla
artificialmente con más reglas empeoraría, no mejoraría, la fiabilidad del
sistema.

\begin{table}[H]
  \centering
  \small
  \renewcommand{\arraystretch}{1.1}
  \begin{tabular}{
    p{0.08\textwidth}
    >{\raggedright\arraybackslash}p{0.66\textwidth}
    p{0.16\textwidth}}
    \hline
    \textbf{Obj.} & \textbf{Descripción resumida} & \textbf{Cumplimiento} \\
    \hline
    OE1 & Despliegue de la infraestructura base de la plataforma & Cumplido \\
    OE2 & Validación del mecanismo central de almacenamiento y procesamiento & Cumplido \\
    OE3 & Incorporación de datos ORCID reales & Cumplido \\
    OE4 & Incorporación de currículos CVN validados frente a Open CVN & Cumplido \\
    OE5 & Integración de ORCID con CVN por medio de Open CVN & Cumplido con garantías parciales \\
    OE6 & Cálculo y publicación consistente de indicadores & Cumplido \\
    OE7 & Visualización de los indicadores en un panel accesible & Cumplido \\
    OE8 & Banco de pruebas de rendimiento por paralelismo y escala & Cumplido \\
    OE9 & Robustez operativa mediante uso real y guía de reproducción & Cumplido \\
    \hline
  \end{tabular}
  \caption{Grado de cumplimiento de los objetivos específicos.}
  \label{tab:tfm_conclusiones_objetivos}
\end{table}

\section{Contribuciones principales}
\label{sec:conclusiones_contribuciones}

Las contribuciones de este trabajo quedan demostradas, con evidencia
concreta, a lo largo del documento:

\begin{itemize}
  \item \textbf{Plataforma \textit{lakehouse} verificada de extremo a
  extremo}: el Capítulo~\ref{cap:infraestructura} demuestra con dos
  comprobaciones independientes que el mecanismo central de
  almacenamiento y procesamiento funciona antes de construir ningún flujo
  de negocio sobre él.
  \item \textbf{Mecanismo de fusión entre ORCID y CVN diseñado desde el
  origen del dato}: el Capítulo~\ref{cap:ingesta} siembra los currículos
  sintéticos con el identificador ORCID de su semilla real desde su
  propia generación, sin enlazar dos fuentes independientes a posteriori.
  \item \textbf{Resolución de identidad determinista con evidencia
  auditable}: el Capítulo~\ref{cap:procesamiento} fusiona registros de
  fuentes distintas mediante reglas explícitas, cada una con su propia
  evidencia, y mide su precisión y su cobertura.
  \item \textbf{Hallazgo del banco de pruebas de rendimiento}: el
  Capítulo~\ref{cap:visualizacion} demuestra que el número de ejecutores
  de Spark es una palanca de fiabilidad, no solo de velocidad.
  \item \textbf{Guía de reproducibilidad validada por una reconstrucción
  real}: la plataforma se reconstruyó desde cero sobre un clúster aislado
  siguiendo únicamente esa guía, en el Capítulo~\ref{cap:visualizacion}.
\end{itemize}

\section{Resultados de aprendizaje desarrollados}
\label{sec:conclusiones_outcomes}

El Capítulo~\ref{cap:introduccion} identificó seis resultados de
aprendizaje del máster trabajados en este proyecto. La
Tabla~\ref{tab:tfm_conclusiones_outcomes} recoge su evidencia concreta.

\begin{table}[H]
  \centering
  \footnotesize
  \renewcommand{\arraystretch}{1.05}
  \begin{tabular}{
    p{0.1\textwidth}
    >{\raggedright\arraybackslash}p{0.8\textwidth}}
    \hline
    \textbf{Resultado} & \textbf{Evidencia concreta} \\
    \hline
    CN02 & Arquitectura \textit{lakehouse} justificada en el
    Capítulo~\ref{cap:estado_arte}, desplegada en el
    Capítulo~\ref{cap:infraestructura}. \\
    HA01 & Adquisición y generación de fuentes enlazadas por diseño en el
    Capítulo~\ref{cap:ingesta}, fusionadas en el
    Capítulo~\ref{cap:procesamiento}. \\
    HA02 & Procesamiento distribuido del Capítulo~\ref{cap:procesamiento},
    evaluado cuantitativamente en el Capítulo~\ref{cap:visualizacion}. \\
    HA03 & Orquestación de flujos ETL de los
    Capítulos~\ref{cap:infraestructura} y~\ref{cap:procesamiento}, con su
    almacenamiento escalable. \\
    CP01 & Planificación y despliegue de una solución Big Data, de los
    Capítulos~\ref{cap:estado_arte},~\ref{cap:infraestructura}
    y~\ref{cap:visualizacion}. \\
    CP04 & Proyecto profesional original de la plataforma completa,
    evaluada en el Capítulo~\ref{cap:visualizacion} y cerrada aquí. \\
    \hline
  \end{tabular}
  \caption{Resultados de aprendizaje del máster y evidencia concreta de cada uno.}
  \label{tab:tfm_conclusiones_outcomes}
\end{table}

\section{Limitaciones y trabajo futuro}
\label{sec:conclusiones_limitaciones}

Los capítulos anteriores han documentado limitaciones concretas sin
ocultarlas, distinguiendo las que proceden de un componente de terceros de
las que son una decisión deliberada de alcance. Esta sección las
consolida junto con el trabajo futuro que se deriva de cada una.

La única de origen externo es el \textbf{fallo conocido del orquestador de
flujos}, sin corrección oficial disponible, documentado con su receta
operativa de recuperación en el Capítulo~\ref{cap:visualizacion}. El
trabajo futuro solo puede consistir en incorporar la corrección oficial
cuando exista.

El resto son decisiones deliberadas de alcance, cada una ya justificada en
su capítulo:

\begin{itemize}
  \sloppy
  \item sin capa de consulta interactiva independiente como
  Trino~\cite{trino_docs}
  \item sin observabilidad centralizada como Prometheus~\cite{prometheus_docs}
  y Grafana~\cite{grafana_docs}
  \item un único nodo local, sin infraestructura como código mediante
  Terraform~\cite{terraform_docs}, en vez de un clúster real en la nube
  \item reconstrucción completa de silver y gold en cada ejecución, en
  lugar de procesamiento incremental
  \item resolución de identidad determinista, frente a las alternativas
  probabilísticas ya citadas en el Capítulo~\ref{cap:procesamiento}
\end{itemize}

La Tabla~\ref{tab:tfm_conclusiones_limitaciones} consolida cada una con su
trabajo futuro asociado.

\begin{table}[H]
  \centering
  \footnotesize
  \renewcommand{\arraystretch}{1.0}
  \begin{tabular}{
    >{\raggedright\arraybackslash}p{0.44\textwidth}
    >{\raggedright\arraybackslash}p{0.44\textwidth}}
    \hline
    \textbf{Limitación} & \textbf{Trabajo futuro asociado} \\
    \hline
    Fallo del orquestador de flujos, sin corrección oficial & Incorporar la corrección oficial cuando esté disponible \\
    Sin capa de consulta interactiva independiente & Incorporar un motor de consulta federado si surge esa necesidad \\
    Sin observabilidad centralizada & Desplegar una pila de observabilidad continua para producción \\
    Único nodo local, sin infraestructura como código para la nube & Automatizar el despliegue en la nube sobre la misma definición \\
    Reconstrucción completa en cada ejecución & Procesar de forma incremental sobre los \textit{snapshots} de Iceberg \\
    Resolución de identidad determinista, sin aprendizaje automático & Explorar vinculación probabilística o basada en aprendizaje automático \\
    \hline
  \end{tabular}
  \caption{Limitaciones documentadas y trabajo futuro asociado.}
  \label{tab:tfm_conclusiones_limitaciones}
\end{table}

\section{Conclusiones}
\label{sec:conclusiones_finales}

El resumen, el cumplimiento de objetivos y las contribuciones ya
presentados permiten cerrar el documento con una valoración conjunta. El
objetivo general se considera cumplido y, con él, los nueve objetivos
específicos, con la única salvedad medida de la cobertura de la
resolución de identidad, una limitación de evidencia y no una afirmación
sin respaldo. Las limitaciones recogidas no comprometen esta conclusión.
Delimitan, con precisión y sin ocultamiento, hasta dónde llega la
propuesta actual. Algunas son ajenas al sistema desarrollado y otras son
decisiones deliberadas de alcance, propias de un Trabajo Fin de Máster y
no de un producto cerrado, y el sistema las documenta en lugar de
disimularlas.

En conjunto, este trabajo demuestra que es posible construir, a partir de
CVN y ORCID, una plataforma \textit{lakehouse} autoalojada que fusiona
información curricular documental con registros públicos a escala masiva,
conserva procedencia verificable en cada capa, y ofrece garantías
reproducibles allí donde el proceso es determinista. Esa plataforma
completa, más que el panel de indicadores o cualquier componente aislado
que la demuestra, es la aportación principal de este Trabajo Fin de
Máster.
